\SourcePage{493}
\SGAsection{13}{与超越理论的联系：复解析情形}
关于\Term{analytic-space}{解析空间}的\Term{generic}{一般}概念，参见 [2]。

\textbf{13.1。$S$ 上\Term{analytic-group}{解析群}的回顾。}
设 $S$ 是一个\Term{analytic-space}{解析空间}，$A$ 是 $S$ 上的一个\Term{analytic-space}{解析空间}，并带有 $S$ 上的\Term{commutative-group}{交换群}结构（即 $A$ 是 $S$ 上\Term{analytic-space}{解析空间}\Term{category}{范畴}中的\Term{commutative-group}{交换群}）。假设 $A$ 在 $S$ 上\Term{smooth}{光滑}。
\SourcePage{494}
记
\begin{equation}
 \underline E=\operatorname{Lie}(A/S)
 \tag{13.1.1}
\end{equation}
为 $S$ 上的\Term{locally-free-module}{局部自由模}，它是 $A$ 在 $S$ 上的\Term{relative-tangent-sheaf}{相对切层}沿\Term{zero-section}{零截面}的\Term{inverse-image}{逆像}。这里我们承认\Term{exponential-homomorphism}{指数同态}
\begin{equation}
 \exp:E\longrightarrow A
 \tag{13.1.2}
\end{equation}
的定义，其中 $E$ 是 $S$ 上的\Term{analytic-vector-bundle}{解析向量丛}，其\Term{holomorphic-section}{全纯截面}\Term{module}{模}为 $\underline E$。在一点 $s\in S$ 的\Term{neighborhood}{邻域}中，先用\Term{characteristic-zero}{特征零}\Term{formal-group}{形式群}理论中熟知的方法，在 $E$ 的\Term{zero-section}{零截面}附近定义这一\Term{homomorphism}{同态}，再显然地利用其\Term{multiplication}{乘法}性质将它延拓到整个 $E$。这是 $S$ 上\Term{relative-analytic-group}{相对解析群}的\Term{homomorphism}{同态}，它是\Term{etale}{平展}的，像是 $A$ 的\Term{open-subspace}{开子空间} $A^0$，即 $A$ 在 $S$ 上各\Term{fiber}{纤维}的\Term{identity-component}{单位分支}之并。若以 $R$ 表示 $\exp$ 的\Term{kernel}{核}所定义的\Term{analytic-subspace}{解析子空间}，便有\Term{relative-analytic-group}{相对解析群}的\Term{exact-sequence}{正合列}
\begin{equation}
 0\longrightarrow R\longrightarrow E\longrightarrow A^0\longrightarrow0,
 \tag{13.1.3}
\end{equation}
其中\Term{morphism}{态射} $E\to A^0$ 是\Term{surjective}{满射}且\Term{etale}{平展}，所以 $R$ 是 $S$ 上的\Term{etale}{平展}\Term{relative-analytic-group}{相对解析群}。它是 $E$ 的一个\Term{analytic-subspace}{解析子空间}；它在 $E$ 中\Term{closed}{闭}，当且仅当 $A^0$ 在 $S$ 上\Term{separated}{分离}，亦即 $A^0$ 的\Term{zero-section}{零截面}在 $A^0$ 中\Term{closed}{闭}。

\SourcePage{495}
\textbf{13.1.4。}当然，\Term{exact-sequence}{正合列} (13.1.3) 关于 $A$ 具有\Term{functorial}{函子性}，并与任意\Term{base-change}{基变换} $S'\to S$ 交换。由此得到一个与\Term{base-change}{基变换}相容的\Term{equivalence-of-categories}{范畴等价}：一方是 $S$ 上\Term{smooth}{光滑}且\Term{fiber}{纤维}\Term{connected}{连通}的\Term{relative-analytic-group}{相对解析群}的\Term{category}{范畴}；另一方是二元组 $(E,R)$ 的\Term{category}{范畴}，其中 $E$ 为 $S$ 上的\Term{locally-free}{局部自由}\Term{vector-bundle}{向量丛}，$R$ 为 $E$ 的\Term{analytic-subgroup}{解析子群}且在 $S$ 上\Term{etale}{平展}。

\textbf{13.1.5。}当 $A$ 是 $S$ 上的“\Term{abeloid-group}{类阿贝尔群}”（abéloïde），即它在 $S$ 上\Term{smooth}{光滑}、\Term{proper}{固有}、\Term{separated}{分离}且\Term{fiber}{纤维}\Term{connected}{连通}，特别地 $A=A^0$ 时，$R$ 是 $S$ 上的\Term{locally-constant-group}{局部常值群}，其在每一点 $s$ 的\Term{rank}{秩}（作为 $\mathbb Z$-\Term{module}{模}的\Term{rank}{秩}）等于 $2\operatorname{rank}_s(\underline E)$；反之亦然。这样便得到了熟知的解释：$S$ 上的\Term{abeloid-space}{类阿贝尔空间}\Term{equivalence}{等价}于 $S$ 上的一个“\Term{local-system}{局部系统}” $R$，以及由 $R$ 定义的实\Term{analytic-bundle}{解析丛} $R\otimes_{\mathbb Z}\mathbb R$ 上的一个复\Term{analytic-vector-bundle}{解析向量丛}结构。

\textbf{13.2。沿一个\Term{fiber}{纤维}的\Term{formal-completion}{形式完备化}。}
从现在起，给定一点 $s\in S$，并在 $s$ 附近作\Term{local-scheme}{局部}讨论。对每个整数 $n\geq0$，以 $S_n$ 表示 $s$ 在 $S$ 中的第 $n$ 个\Term{infinitesimal-neighborhood}{无穷小邻域}；对 $S$ 上任一\Term{analytic-space}{解析空间} $X$，同样令 $X_n=X\times_S S_n$。当 $n$ 变化时，$X_n$ 构成 $S$ 上\Term{analytic-space}{解析空间}的一个\Term{inductive-system}{归纳系统}，记为 $\widehat X$，称为 $X$ 沿 $X_s$ 的\Term{formal-completion}{形式完备化}。显然，它关于 $X$ 具有\Term{functorial}{函子性}，且其形成与\Term{finite}{有限}\Term{projective-limit}{射影极限}交换；
\SourcePage{496}
因而 $X\mapsto\widehat X$ 将群变为群。还应注意，$\widehat X$ 实际上只依赖于 $X$ 在 $S$ 的任意一个\Term{open-neighborhood}{开邻域} $U$ 上的\Term{restriction}{限制}。

\begin{SGAresult}{\Term{lemma}{引理}}{13.2.1}
数据如 13.1，并假设满足下列条件：a) $A=A^0$；b) $R$ 在 $s$ 附近是\Term{constant-group}{常值群}；c) $R_s$ 生成复\Term{vector-space}{向量空间} $E_s$。设 $A'$ 是 $S$ 上另一个\Term{smooth}{光滑}\Term{analytic-group}{解析群}，考虑映射
\[
 u\longmapsto\widehat u:
 \operatorname{Hom}_{S\text{-gr}}(A,A')\longrightarrow
 \operatorname{Hom}_{\widehat S\text{-gr}}(\widehat A,\widehat A'),
\]
其中记号同 13.2。再考虑将它\Term{restriction}{限制}在 $s$ 的变化着的\Term{open-neighborhood}{开邻域} $U$ 上，并取\Term{inductive-limit}{归纳极限}而得到的类似映射
\begin{equation}
 u\longmapsto\widehat u:
 \operatorname{Hom}_{(S,s)\text{-gr}}(A,A')\longrightarrow
 \operatorname{Hom}_{\widehat S\text{-gr}}(\widehat A,\widehat A'),
 \tag{13.2.1.1}
\end{equation}
其中 $(S,s)$ 表示 $S$ 在 $s$ 处的\Term{analytic-space-germ}{解析空间芽}。则后一映射是\Term{bijective}{双射}。
\end{SGAresult}
\begin{proof}
考虑 $A'$ 所定义的、与 (13.1.3) 类似的\Term{exact-sequence}{正合列}
\[
 0\longrightarrow R'\longrightarrow E'\longrightarrow A'^0\longrightarrow C.
\]
因为\Term{homomorphism}{同态} $A\to A'$ 必经 $A'^0$ \Term{resolution}{分解}，而且在作\Term{base-change}{基变换} $U\to S$ 和 $S_n\to S$ 后也一样，所以可以假设 $A'=A'^0$。由 13.1.4 的对应，只须证明：给定一个将 $R$ 映入 $R'$ 的\Term{vector-bundle}{向量丛}\Term{homomorphism}{同态} $E\to E'$ 的芽，
\SourcePage{497}
\Term{equivalence}{等价}于给定一个将 $R_n$ 映入 $R'_n$ 的\Term{homomorphism}{同态} $E_n\to E'_n$ 的相容系统。若 $\mathcal O_s$ 为 $S$ 在 $s$ 处的\Term{local-ring}{局部环}，则给定这样一个相容的\Term{homomorphism}{同态}系统 $E_n\to E'_n$，\Term{equivalence}{等价}于给定一个 $\widehat{\mathcal O}_s$-\Term{module}{模}\Term{homomorphism}{同态}
\[
 \widehat F:\widehat{\underline E}_s\longrightarrow\widehat{\underline E}'_s
\]
（$\underline E_s$ 和 $\underline E'_s$ 是对应于 $E,E'$ 的 $\mathcal O_S$-\Term{module}{模} $\underline E,\underline E'$ 在 $s$ 处的\Term{stalk}{茎}，它们是\Term{free-module}{自由} $\mathcal O_s$-\Term{module}{模}）；而给定\Term{vector-bundle}{向量丛}\Term{homomorphism}{同态} $E\to E'$ 的一个芽，则\Term{equivalence}{等价}于给定一个 $\mathcal O_s$-\Term{module}{模}\Term{homomorphism}{同态}
\[
 F:\underline E_s\longrightarrow\underline E'_s.
\]
这已说明映射 (13.2.1.1) 是\Term{injective}{单射}。至于\Term{surjective}{满射}性，只须证明：若给定上述\Term{homomorphism}{同态} $\widehat F$，且对每个 $n$，它诱导的\Term{homomorphism}{同态} $F_n:E_n\to E'_n$ 将 $R_n$ 映入 $R'_n$，则 $\widehat F$ 来自一个 $F$，即将 $\underline E_s$ 映入 $\underline E'_s$，且相应的\Term{homomorphism}{同态}芽 $E\to E'$ 将 $R$ 映入 $R'$。由假设 c)，$R_s=R_0$ 生成 $E_s$，故由 \Term{nakayama-lemma}{Nakayama 引理}，群 $R_0$（将它与 $R$ 在 $S$ 上的\Term{section}{截面}群认同，或与 $R$ 在 $s$ 附近的\Term{section-germ}{截面芽}群认同，均可）生成 $\underline E_s$。所以，为了证明 $\widehat F(\underline E_s)\subset\underline E'_s$，只须证明 $\widehat F(R_0)\subset\underline E'_s$。但 $\widehat F(R_0)\subset R'_0$，这里将 $R'_0$ 本身解释为 $R'\subset E'$ 在 $s$ 附近的\Term{section-germ}{截面芽}群，因而得到所需包含关系，并得到 $F$ 的存在性。此外可见 $F(R_0)\subset R'_0$，这里 $R_0,R'_0$ 分别视为 $E,E'$ 的\Term{section-germ}{截面芽}群。根据 b)，$R_0$ 在 $s$ 附近生成 $R$，
\SourcePage{498}
故 $F(R)\subset R'$ 在 $s$ 附近成立，证明完毕。
\end{proof}

\begin{SGAresult}{推论}{13.2.2}
若 $A$ 满足 13.2.2 的条件 a)、b)、c)，则由 $\widehat S$-群 $\widehat A$ 可在唯一\Term{isomorphism}{同构}意义下确定 $A$ 在 $s$ 附近的情形。
\end{SGAresult}

\textbf{13.3。解析的 \Term{raynaud-extension}{Raynaud 扩张}。}
现在可以在\Term{analytic-geometry}{解析几何}中重述第 7 节给出的\Term{raynaud-extension}{Raynaud 扩张}的构造。设给定如 13.1 中的 $A$，并有 $A=A_0$；考虑 $R_0=R_s$，它是一个\Term{finite-type}{有限型}\Term{free-module}{自由} $\mathbb Z$-\Term{module}{模}。将 $S$ 缩小为 $s$ 的一个足够小的\Term{open-neighborhood}{开邻域}，并将 $R_0$ 解释为 $R$ 在 $s$ 附近的\Term{section-germ}{截面芽}群，便得到一个\Term{homomorphism}{同态}
\begin{equation}
 R_{0S}\longrightarrow R\subset E,
 \tag{13.3.1}
\end{equation}
其中 $R_{0S}$ 为 $S$ 上取值为 $R_0$ 的\Term{constant-analytic-group}{常值解析群}。将 $s$ 的\Term{neighborhood}{邻域}取得足够小，使得 $\mathcal O_S$ 的任一在 $s$ 处的芽为零的\Term{section}{截面}都为零，即可看出上述\Term{homomorphism}{同态}是\Term{injective}{单射}。因此其像是 $R$ 的一个\Term{constant-analytic-subgroup}{常值解析子群}；显然，它是 $R$ 的最大\Term{constant-subgroup}{常值子群}，称为 $R$ 的\Term{fixed-part}{固定部分}，记为 $R^f$：
\begin{equation}
 R_{0S}\xrightarrow{\sim}R^f\lhook\joinrel\longrightarrow R,
 \qquad\text{诱导}\quad R^f_s\xrightarrow{\sim}R_s.
 \tag{13.3.2}
\end{equation}
于是可以考虑\Term{quotient-group}{商群}
\begin{equation}
 A^\natural=E/R^f.
 \tag{13.3.3}
\end{equation}
\SourcePage{499}
由 (13.1.3)，得到一个\Term{canonical}{典范}的\Term{exact-sequence}{正合列}
\begin{equation}
 0\longrightarrow R/R^f\xrightarrow{i}A^\natural\longrightarrow A\longrightarrow0.
 \tag{13.3.4}
\end{equation}
注意，只须将 $S$ \Term{restriction}{限制}在 $s$ 的一个\Term{neighborhood}{邻域}内，$R^f$ 在 $E$ 中总是\Term{closed}{闭}的，因而 $A^\natural$ 在 $S$ 上\Term{separated}{分离}。还应注意，\Term{homomorphism}{同态} $A^\natural\to A$ 是\Term{etale}{平展}的，故诱导\Term{isomorphism}{同构}
\begin{equation}
 \operatorname{Lie}(A^\natural)\xrightarrow{\sim}\operatorname{Lie}(A)=\underline E.
 \tag{13.3.5}
\end{equation}
另一方面，对每个 $n\geq0$，\Term{canonical}{典范}\Term{homomorphism}{同态} $A^\natural\to A$ 诱导\Term{isomorphism}{同构} $A^\natural_n\xrightarrow{\sim}A_n$（因为 $R^f_n\xrightarrow{\sim}R_n$），所以它诱导一个\Term{isomorphism}{同构}
\begin{equation}
 \widehat{A^\natural}\xrightarrow{\sim}\widehat A.
 \tag{13.3.6}
\end{equation}
后一事实表明，$A^\natural$ 确实起着我们在第 7 节考察过、同样记为 $A^\natural$ 的群的\Term{action}{作用}。这里构造它时，并未对\Term{fiber}{纤维} $A_s=A_0$ 作“由\Term{torus}{环面}\Term{extension}{扩张}一个\Term{abeloid-group}{类阿贝尔群}”之类的假设；不过，要在 $A^\natural$ 上恢复重要的\Term{extension}{扩张}结构 (7.2.3)，这个假设是必需的。另一方面，(13.3.4) 的\Term{homomorphism}{同态} $A^\natural\to A$ 本质上是\Term{transcendental}{超越}性的，在第 7 节的纯\Term{algebra}{代数}理论中没有对应物。至于 (13.3.4) 中\Term{homomorphism}{同态} $i:R/R^f\to A^\natural$ 的纯\Term{algebra}{代数}解释，我们将在第 14 节再次讨论。

先假设给定一个\Term{complex-subtorus}{复子环面}
\begin{equation}
 T_0\subset A_0,
 \tag{13.3.7}
\end{equation}
这里“\Term{complex-torus}{复环面}”是指\Term{isomorphism}{同构}于 $\mathbb C^*$ 的某个\Term{cartesian}{笛卡尔}幂的\Term{analytic-group}{解析群}。
\SourcePage{500}
令
\begin{equation}
 T=T_0\times S
 \tag{13.3.8}
\end{equation}
为 $S$ 上取值为 $T_0$ 的\Term{constant-torus}{常值环面}。仿照 5.1 的做法，得到\Term{analytic-group}{解析群}的\Term{inductive-system}{归纳系统}之间唯一的\Term{homomorphism}{同态}（甚至是\Term{immersion}{浸入}）
\[
 \widehat T\longrightarrow\widehat{A^\natural},
\]
它在 $S_0$ 上化为 (13.3.7)。根据 13.2.2，这个\Term{homomorphism}{同态}来自唯一的\Term{homomorphism}{同态}芽
\begin{equation}
 T\longrightarrow A^\natural.
 \tag{13.3.9}
\end{equation}
将 $S$ 在 $s$ 附近取得足够小，便可假设这个\Term{homomorphism}{同态}在整个 $S$ 上有定义。相应的\Term{homomorphism}{同态}
\[
 \operatorname{Lie}(T)\longrightarrow\operatorname{Lie}(A^\natural)\simeq\underline E
\]
在各 $S_n$ 上是\Term{injective}{单射}，所以在 $\widehat S$ 上是\Term{injective}{单射}，进而由\Term{flat-descent}{平坦下降}，在 $\mathcal O_s$ 上也是\Term{injective}{单射}。故存在 $S$ 的一个\Term{neighborhood}{邻域}，可以将它取为 $S$ 本身，使 $\operatorname{Lie}(T)$ 与 $\underline E$ 的一个\Term{locally-direct-summand}{局部直和因子}\Term{submodule}{子模}相认同；将它记为 $\underline E^t$，并以 $E^t$ 表示 $E$ 的相应\Term{subbundle}{子丛}：
\begin{equation}
 \operatorname{Lie}(T)\xrightarrow{\sim}\underline E^t\subset\underline E
 \xrightarrow{\sim}\operatorname{Lie}(A^\natural)
 \xrightarrow{\sim}\operatorname{Lie}(A),\qquad E^t\subset E.
 \tag{13.3.10}
\end{equation}
同样，以
\begin{equation}
 R_0^t\subset R_0,
 \qquad R^t=R_{0S}^t\subset R_{0S}\xrightarrow{\sim}R^f
 \tag{13.3.11}
\end{equation}
表示由 $A_0$ 的\Term{subtorus}{子环面} $T_0$ 所定义的 $R_0$ 的\Term{sublattice}{子格}，以及相应的 $S$ 上\Term{constant-analytic-group}{常值解析群}。这样，\Term{homomorphism}{同态} (13.3.9)
\SourcePage{501}
就是由包含关系
\[
 E^t\lhook\joinrel\longrightarrow E,\qquad
 R^t\lhook\joinrel\longrightarrow R^f
\]
诱导的\Term{homomorphism}{同态}
\[
 E^t/R^t_{0S}=E^t/R^t\longrightarrow E/R_{0S}=E/R^f.
\]
利用 (13.3.9) 在 $s$ 处的\Term{fiber}{纤维}上诱导一个\Term{immersion}{浸入}这一事实，可以得出\Term{homomorphism}{同态}
\[
 R_0/R_0^t\longrightarrow E_0/E_0^t=(E/E^t)_0
\]
是\Term{injective}{单射}。由此立刻得到，对邻近 $s$ 的 $s'$，
\[
 R_0/R_0^t\longrightarrow(E/E^t)_{s'}
\]
仍为\Term{injective}{单射}。这意味着，在 $s$ 附近缩小 $S$ 后，可以假设 (13.3.9) 是一个\Term{closed-immersion}{闭浸入}。

于是自然应引入
\begin{equation}
 B=A^\natural/T\simeq(E/E^t)/(R^f/R^t),
 \tag{13.3.12}
\end{equation}
从而在 $A^\natural$ 上得到\Term{extension}{扩张}结构
\begin{equation}
 0\longrightarrow T\longrightarrow A^\natural\longrightarrow B\longrightarrow0,
 \tag{13.3.13}
\end{equation}
其中 $B$ 是 $S$ 上的\Term{relative-analytic-group}{相对解析群}，在 $S$ 上\Term{smooth}{光滑}、\Term{separated}{分离}且\Term{fiber}{纤维}\Term{connected}{连通}。由构造本身，有
\begin{equation}
 \operatorname{Lie}(B)\simeq\underline E/\underline E^t,
 \tag{13.3.14}
\end{equation}
以及 $B$ 的\Term{exponential-exact-sequence}{指数正合列}
\SourcePage{502}
\begin{equation}
 0\longrightarrow R^f/R^t\longrightarrow E/E^t\longrightarrow B\longrightarrow0.
 \tag{13.3.15}
\end{equation}
注意，
\begin{equation}
 R^f/R^t=(R_0/R_0^t)_S
 \tag{13.3.16}
\end{equation}
是 $S$ 上的\Term{constant-analytic-group}{常值解析群}。所以由 13.1.5，$B$ 在 $s$ 的一个\Term{neighborhood}{邻域}上是 $S$ 上的\Term{abeloid-group}{类阿贝尔群}，当且仅当它在 $s$ 处的\Term{fiber}{纤维}
\[
 B_0=A_0/T_0
\]
是一个\Term{abeloid-analytic-group}{类阿贝尔解析群}；此时 $A_0$ 就成为由\Term{torus}{环面} $T_0$ \Term{extension}{扩张}这样的群 $B_0$ 所得到的群。将 $A_0$ 写成由\Term{complex-torus}{复环面}\Term{extension}{扩张}一个\Term{abeloid-group}{类阿贝尔群}的\Term{form}{形式} $(*)$，每一种这样的写法，都在 $s$ 附近的 $A^\natural$ 上给出唯一的\Term{extension}{扩张}结构 (13.3.13)，即由\Term{torus}{环面} $T$ \Term{extension}{扩张}一个\Term{abeloid-group}{类阿贝尔群} $B$，并在 $s$ 处诱导给定的\Term{extension}{扩张}结构。

\textbf{13.4。用 Deligne 的“\Term{mixed-hodge-structure}{混合 Hodge 结构}”作解释。}
先回到 13.1.4 的对应。我们将稍作改述，使它更适合引入此处所需的“\Term{hodge-structure}{Hodge 结构}”。为此，以 $\underline R$ 表示 $R$ 在 $S$ 上的\Term{section}{截面}\Term{sheaf}{层}，并引入 $\mathcal O_S$-\Term{module}{模}
\begin{equation}
 \mathscr E(A)=\mathscr E=\underline R\otimes_{\mathbb Z}\mathcal O_S;
 \tag{13.4.1}
\end{equation}
当 $R$ \Term{locally-constant}{局部常值}时，它是\Term{finite-type}{有限型}\Term{locally-free}{局部自由}的。（注意，若 $R$ 并非\Term{locally-constant}{局部常值}，则\Term{sheaf}{层} $\mathscr E$ 未必\Term{coherent}{凝聚}。）\Term{homomorphism}{同态} (13.1.3) $R\to E$ 定义一个\Term{canonical}{典范}\Term{homomorphism}{同态}
\SourcePage{503}
\begin{equation}
 \mathscr E=\underline R\otimes_{\mathbb Z}\mathcal O_S\longrightarrow\underline E,
 \tag{13.4.2}
\end{equation}
其\Term{kernel}{核}记为 $\underline F(A)$ 或简记为 $\underline F$，从而有\Term{exact-sequence}{正合列}
\begin{equation}
 0\longrightarrow\underline F\longrightarrow\mathscr E\longrightarrow\underline E.
 \tag{13.4.3}
\end{equation}
最有意义的情形，是\Term{homomorphism}{同态} (13.4.2) 为\Term{epimorphism}{满态射}，即对每个 $s\in S$，$R_s$ 生成复\Term{vector-space}{向量空间} $E_s$。此时便得到\Term{exact-sequence}{正合列}
\begin{equation}
 0\longrightarrow\underline F\longrightarrow\mathscr E\longrightarrow\underline E\longrightarrow0;
 \tag{13.4.4}
\end{equation}
若 $R$ \Term{locally-constant}{局部常值}，这就是 $S$ 上\Term{locally-free-module}{局部自由模}的一个\Term{short-exact-sequence}{短正合列}。此时，知道二元组 $(E,R)$ \Term{equivalence}{等价}于知道二元组 $(\underline R,\underline F)$，其中 $\underline R$ 是 $S$ 上的\Term{locally-constant}{局部常值}\Term{analytic-group}{解析群}，其\Term{fiber}{纤维}为\Term{finite-type}{有限型}\Term{free-module}{自由} $\mathbb Z$-\Term{module}{模}，而 $\underline F$ 是 $\mathscr E=\underline R\otimes_{\mathbb Z}\mathcal O_S$ 的一个\Term{locally-direct-summand}{局部直和因子}\Term{submodule}{子模}，此外还满足关系
\begin{equation}
 \underline F\cap\underline R=0.
 \tag{13.4.5}
\end{equation}
从\Term{hodge-structure}{Hodge 结构}的观点，应把 $\mathscr E$ 看成带有“\Term{hodge-filtration}{Hodge 滤过}”
\begin{equation}
 \operatorname{Fil}^{1}(\mathscr E)=\mathscr E,\qquad
 \operatorname{Fil}^{0}(\mathscr E)=\underline F,\qquad
 \operatorname{Fil}^{-1}(\mathscr E)=0.
 \tag{13.4.6}
\end{equation}
注意，\Term{exact-sequence}{正合列} (13.4.3) 仍然\Term{functorial}{函子性}地依赖于 $S$ 上的\Term{analytic-group}{解析群} $A$，且与任意\Term{base-change}{基变换}交换。

\SourcePage{504}
现在回到 13.3 的条件，令
\begin{equation}
 \mathscr E^f=\underline R^f\otimes_{\mathbb Z}\mathcal O_S,\qquad
 \mathscr E^t=\underline R^t\otimes_{\mathbb Z}\mathcal O_S.
 \tag{13.4.7}
\end{equation}
这些\Term{module}{模}定义在 $s$ 的某个适当\Term{neighborhood}{邻域}上（可将它取为 $S$ 本身），从而定义 $\mathscr E$ 的一个\Term{filtration}{滤过}（称为\Term{weight-filtration}{权滤过}）
\begin{equation}
 \mathscr E\supset\mathscr E^f\supset\mathscr E^t\supset0;
 \tag{13.4.8}
\end{equation}
这里应将 $\mathscr E^t,\mathscr E^f,\mathscr E$ 分别看成\Term{weight}{权}为 $-2,-1,0$ 的部分（所以这是一个\Term{increasing-filtration}{递增滤过}）。注意，$\mathscr E^t$ 和 $R^t$ 一样，依赖于 $A_0=A_s$ 的\Term{subtorus}{子环面} $T_0$ 的选择。还应注意，因为 $R^f$ 和 $R^t$ \Term{locally-constant}{局部常值}，$\mathscr E^f$ 和 $\mathscr E^t$ 是\Term{locally-free}{局部自由}的，且也\Term{solvable}{可解}释为
\begin{equation}
 \mathscr E^f=\mathscr E(A^\natural),\qquad
 \mathscr E^t=\mathscr E(T).
 \tag{13.4.9}
\end{equation}
将\Term{exact-sequence}{正合列} (13.4.3) 关于 $A$ 的\Term{functoriality}{函子性}用于\Term{homomorphism}{同态}
\[
 T\longrightarrow A^\natural\longrightarrow A,
\]
立刻可知 $\underline F(A^\natural)$ 和 $\underline F(T)$ 分别认同于
\begin{equation}
 \underline F(A^\natural)\simeq\underline F^f
 \mathrel{\mathop:}=\underline F\cap\mathscr E^f,
 \qquad
 \underline F(T)=\underline F^t=\underline F\cap\mathscr E^t=0.
 \tag{13.4.10}
\end{equation}
这里 $\underline F^t=0$ 来自显然的事实：对\Term{torus}{环面} $T$，有 $\underline F(T)=0$，而 (13.4.3) 型的\Term{exact-sequence}{正合列}在此变为
\[
 0\longrightarrow0\longrightarrow\mathscr E(T)\longrightarrow\operatorname{Lie}(T)\longrightarrow0.
\]
同时可见，复合\Term{homomorphism}{同态}
\[
 \mathscr E^t\longrightarrow\mathscr E\longrightarrow\underline E
\]
诱导一个\Term{isomorphism}{同构}
\SourcePage{505}
\begin{equation}
 \mathscr E^t\xrightarrow{\sim}\underline E^t\quad(\simeq\operatorname{Lie}(T)).
 \tag{13.4.11}
\end{equation}
现在假设 $R_s$ 生成复\Term{vector-space}{向量空间} $E_s$。特别地，在 13.3.16 所述的有利情形，即 $A_s$ 是由\Term{complex-torus}{复环面} $T_0$ \Term{extension}{扩张}一个\Term{abeloid-group}{类阿贝尔群}而得到的群时，就满足这一条件。于是，将 $S$ \Term{restriction}{限制}在 $s$ 的某个适当\Term{neighborhood}{邻域}后，(13.4.2) 型的\Term{canonical}{典范}\Term{homomorphism}{同态}（将 $A$ 换为 $A^\natural$）
\[
 \mathscr E^f=\mathscr E(A^\natural)\longrightarrow\operatorname{Lie}(A^\natural)\simeq\underline E
\]
是\Term{epimorphism}{满态射}，因此 (13.4.2) 更是\Term{epimorphism}{满态射}。所以有\Term{exact-sequence}{正合列} (13.4.4) 以及
\begin{equation}
 0\longrightarrow\underline F^f\longrightarrow\mathscr E^f\longrightarrow\underline E\longrightarrow0.
 \tag{13.4.12}
\end{equation}
特别地，因为 $\mathscr E/\underline F\simeq\underline E$，便有
\begin{equation}
 \underline F+\mathscr E^f=\mathscr E.
 \tag{13.4.13}
\end{equation}
此外，因为 $B=A^\natural/T$ 是 $A$ 的一个商，它满足类似的\Term{surjectivity}{满射性}条件，亦即产生一个 (13.4.4) 型的\Term{exact-sequence}{正合列}：
\begin{equation}
 0\longrightarrow\underline F(B)\longrightarrow
 \text{\SGAsourcegap{(13.4.14)中间项在(B)前的记号原缺失}}(B)
 \longrightarrow\operatorname{Lie}(B)\longrightarrow0.
 \tag{13.4.14}
\end{equation}
它用已经引入的\Term{module}{模}便可直接解释。事实上，由 (13.3.12)，
\begin{equation}
 \left\{\begin{aligned}
 \mathscr E(B)&\simeq(\underline R^f/\underline R^t)\otimes_{\mathbb Z}\mathcal O_S
              \simeq\mathscr E^f/\mathscr E^t,\\
 \operatorname{Lie}(B)&\simeq\underline E/\underline E^t.
 \end{aligned}\right.
 \tag{13.4.15}
\end{equation}
\SourcePage{506}
所以，由 (13.4.11) 以及关系 $\underline F\cap\mathscr E^t=0$（13.4.10），可知
\[
 \underline F(B)=\ker(\mathscr E^f/\mathscr E^t\longrightarrow\underline E/\underline E^t)
\]
由\Term{homomorphism}{同态}
\begin{equation}
 \underline F^f\xrightarrow{\sim}\underline F(A^\natural)
 \xrightarrow{\sim}\underline F(B)
 \tag{13.4.16}
\end{equation}
与 $\underline F^f$ \Term{isomorphism}{同构}。

现在汇总“\Term{weight-filtration}{权滤过}” (13.4.8) 的上述性质，以及它与 $\mathscr E$ 上由\Term{submodule}{子模} $\underline F$ 定义的“\Term{hodge-filtration}{Hodge 滤过}” (13.4.6) 的关系。我们取 13.3.16 所述的有利情形，即 $A_s$ 是由\Term{torus}{环面} $T_0$ \Term{extension}{扩张}一个\Term{abeloid-analytic-group}{类阿贝尔解析群} $B_0$ 得到的群（用 $T_0$ 来定义 $R^t$，从而 $B$ 是 $S$ 上的\Term{abeloid-analytic-group}{类阿贝尔解析群}）：

\begin{SGAresult}{命题}{13.4.17}
在上述条件下，$\mathscr E=\underline R\otimes_{\mathbb Z}\mathcal O_S$ 上的“\Term{hodge-filtration}{Hodge 滤过}”和“\Term{weight-filtration}{权滤过}”
\begin{equation}
 \mathscr E\supset\underline F,\qquad
 \mathscr E\supset\mathscr E^f\supset\mathscr E^t\supset0
 \tag{13.4.17.1}
\end{equation}
满足下列条件：

\textup{a)} \Term{weight-filtration}{权滤过}由 $\underline R$ 上的类似\Term{filtration}{滤过}
\[
 \underline R\supset\underline R^f\supset\underline R^t\supset0
\]
经\Term{tensor-product}{张量积} $\otimes_{\mathbb Z}\mathcal O_S$ 得到，其中 $R^f$ 和 $R^t$ 是 $S$ 上的\Term{locally-constant}{局部常值}\Term{analytic-group}{解析群}，取值于\Term{finite-type}{有限型}\Term{free-module}{自由} $\mathbb Z$-\Term{module}{模}（即 $R^f_0=R_s=\pi_1(A_s)$ 和 $R^t_0=\pi_1(T_0)$）。

\textup{b)} 有关系
\begin{equation}
 \underline F\cap\mathscr E^t=0,\qquad
 \underline F+\mathscr E^f=\mathscr E.
 \tag{*}
\end{equation}
\SourcePage{507}
换言之，$\mathscr E^t$ 上由\Term{hodge-filtration}{Hodge 滤过}所诱导的\Term{filtration}{滤过}（采用 (13.4.6) 中的\Term{degree}{次数}约定）集中在\Term{degree}{次数} $-1$；$\mathscr E/\mathscr E^f$ 上由\Term{hodge-filtration}{Hodge 滤过}所诱导的\Term{filtration}{滤过}集中在\Term{degree}{次数} $0$。

\textup{c)} 考虑
\[
 \mathscr E^{\mathrm{ab}}\mathrel{\mathop:}=\mathscr E^f/\mathscr E^t
 \simeq\underline R^{\mathrm{ab}}\otimes_{\mathbb Z}\mathcal O_S,
 \qquad\underline R^{\mathrm{ab}}=\underline R^f/\underline R^t
 \quad\text{（参见 a)）}
\]
上由 $\mathscr E$ 的\Term{hodge-filtration}{Hodge 滤过}诱导的\Term{filtration}{滤过}；它由 $\mathscr E^{\mathrm{ab}}$ 的\Term{submodule}{子模} $\underline F^{\mathrm{ab}}$ 定义，后者是
\[
 \underline F^f\mathrel{\mathop:}=\underline F\cap\mathscr E^f
 \longrightarrow\mathscr E^{\mathrm{ab}}
\]
的像（由第一个关系 $(*)$，此\Term{arrow}{箭头}为\Term{monomorphism}{单态射}）。则二元组
\[
 (\underline R^{\mathrm{ab}},
 \underline F^{\mathrm{ab}}\subset
 \underline R^{\mathrm{ab}}\otimes_{\mathbb Z}\mathcal O_S=\mathscr E^{\mathrm{ab}})
\]
与按照 (13.4.4) 之后的说明、由 $S$ 上某个\Term{abeloid-group}{类阿贝尔群}关联得到的二元组 $(\underline R(B),\underline F(B))$ \Term{isomorphism}{同构}（因而 $B\simeq\mathscr E^{\mathrm{ab}}/(F^{\mathrm{ab}}+R^{\mathrm{ab}})$，其中不带下划线的\Term{signature}{符号}表示相应的 $S$ 上\Term{analytic-group}{解析群}）。
\end{SGAresult}

\textbf{注记 13.4.18。}采用 Deligne [4] 的术语，将前一命题的结论\Term{restriction}{限制}在 $S$ 的一个使 $R$ \Term{locally-constant}{局部常值}的\Term{open-subset}{开子集} $U$ 上，它便表示：$R$ 连同 $\mathscr E=\underline R\otimes_{\mathbb Z}\mathcal O_S$ 上的两个\Term{filtration}{滤过} (13.4.17.1)，是“$U$ 上\Term{mixed-hodge-structure}{混合 Hodge 结构}的一个解析族”。刚才阐明的情形，是 Deligne 在 [9, 9.8 至 9.13] 中提出的高维 Néron 理论猜想的原型。不过，当假设 $A$ 在 $S$ 中某个\Term{nowhere-dense}{稀疏}\Term{analytic-subset}{解析子集} $Z$ 的补\Term{open-subset}{开子集} $U$ 上是\Term{abeloid-group}{类阿贝尔群}时（$\dim S=1$ 时，不妨取 $Z=\{s\}$），
\SourcePage{508}
第 7 节所研究的\Term{algebra}{代数}理论中有一个重要特点未包含在 13.4.17 的描述中。此时由 13.1.5，$R|U$ 是 $U$ 上的\Term{local-system}{局部系统}；所欠缺的是：$U$ 上的\Term{local-system}{局部系统} $R/R^t$ 应当是\Term{constant}{常值}的，更确切地说，它应当与\Term{local-system}{局部系统} $R'^f$ \Term{dual}{对偶}。这里 $R'$ 是 $S$ 上的\Term{etale}{平展}\Term{analytic-group}{解析群}，它与 $S$ 上某个适当的\Term{analytic-group}{解析群} $A'$ 相对应，如同 $R$ 与 $A$ 相对应一样，并使 $A'|U$ 成为 $A|U$ 的“\Term{dual}{对偶}”\Term{abeloid-group}{类阿贝尔群}。（用第 5 节的记号，此断言的\Term{algebraic-correspondence}{代数对应}是由\Term{duality-theorem}{对偶定理} 5.2 得出的 $T_\ell(A_K)/T_\ell(A_K)^t$ 与 $T_\ell(A'_K)^f$ 的\Term{duality}{对偶性}（5.5.9）。）我不知道仅在 13.4.17 的假设下，这一断言是否成立，甚至不知道 $\pi_1(U,t)$ 在 $R_t$ 上的\Term{action}{作用}是否\Term{unipotent}{幺幂}（$t\in U$）。看来，即使再假设 $A|U$ 带有一个“\Term{polarization}{极化}”，即一个在每条\Term{fiber}{纤维}上都\Term{ample}{丰沛}的\Term{invertible-sheaf}{可逆层}，本节使用的初等方法似乎也不能证明这样的结果。不过，研究“\Term{semisimple-group}{半单群}的那些人”（semi-simple group people）似乎知道如何证明，至少在 $S$ \Term{nonsingular}{非奇异}且\Term{dimension}{维数}为 $1$ 时如此，参见报告 [9]。在下一小节中，我们将通过对 $A$ 作相对\Term{algebra}{代数}性的假设，说明一切都运作良好，同时建立与第 7 节构造之间的联系。

\SourcePage{509}
\textbf{13.5。相对\Term{algebra}{代数}的情形。}
现在假设 $A$ 来自 $S$ 上的一个\Term{relative-group-scheme}{相对群概形}，它\Term{smooth}{光滑}、\Term{separated}{分离}，且各\Term{fiber}{纤维}都是由\Term{torus}{环面}\Term{extension}{扩张}\Term{abelian-scheme}{阿贝尔概形}而得到的群。仍以字母 $A$ 表示这个\Term{relative-scheme}{相对概形}。（允许这样滥用记号，是因为\Term{functor}{函子} $G\mapsto G^{\mathrm{an}}$，从所考察类型的\Term{relative-group-scheme}{相对群概形}的\Term{category}{范畴}到 $S$ 上\Term{analytic-group}{解析群}的\Term{category}{范畴}，是\Term{fully-faithful}{全忠实}的；用相对的 GAGA 型结果 [10] 很容易证明这一点。）关于\Term{relative-scheme}{相对概形}的概念，以及 $S$ 上\Term{relative-scheme}{相对概形}与 $S$ 上\Term{analytic-space}{解析空间}之间的关系，均参见 [10]（另参见 [2, 第 15 讲]）。若 $\mathcal O$ 表示 $S$ 在 $s$ 处的\Term{local-ring}{局部环}，则\Term{relative-scheme}{相对概形} $A$ 在 $s$ 附近由它在 $s$ 处\Term{localization}{局部化}得到的、$\operatorname{Spec}(\mathcal O)$ 上的通常\Term{scheme}{概形} $A_{\mathcal O}$ 确定。借助上述记号上的滥用，将 $A_n=A\times_S S_n$ 与 $S_n$ 上的\Term{algebraic-scheme}{代数概形}认同后，显然 $A_{\mathcal O}$ 沿其\Term{special-fiber}{特殊纤维}（\Term{isomorphism}{同构}于 $A_s=A_0$）的\Term{formal-completion}{形式完备化}，与\Term{analytic-group}{解析群} $A$ 沿 $A_s$ 的\Term{formal-completion}{形式完备化}（13.2）\Term{isomorphism}{同构}。在有利的情形，即将 $\widehat A$ 看成 $\widehat{\mathcal O}$ 上的\Term{formal-group-scheme}{形式群概形}时它\Term{algebraizable}{可代数化}（7.2.1）时，便得到 $\operatorname{Spec}(\widehat{\mathcal O})$ 上的一个\Term{group-scheme}{群概形} $A^\natural_{\widehat{\mathcal O}}$，它是由\Term{torus}{环面}\Term{extension}{扩张}一个\Term{abelian-scheme}{阿贝尔概形} $B_{\widehat{\mathcal O}}$ 而得到的群。特别地，若 $S$ 在 $s$ 处\Term{normal}{正规}，且 $A$ 在 $S$ 上\Term{quasi-projective}{拟射影}，便属于这种情形（7.1.5）。为简单起见，假设 $B_{\widehat{\mathcal O}}$ 在 $\widehat{\mathcal O}$ 上\Term{projective}{射影}（已知若 $\mathcal O$，因而 $\widehat{\mathcal O}$，\Term{normal}{正规}，则确实如此）。下面证明：

\SourcePage{510}
\begin{SGAresult}{命题}{13.5.1}
在上述条件下，(13.3.3) 的\Term{analytic-group}{解析群} $A^\natural$ 在 $s$ 附近来自一个在唯一\Term{isomorphism}{同构}意义下确定的\Term{relative-group-scheme}{相对群概形}，它是由\Term{torus}{环面}\Term{extension}{扩张}一个 $S$ 上的\Term{projective}{射影}\Term{relative-abelian-scheme}{相对阿贝尔概形}而得到的群。若 $A^\natural_{\mathcal O}$ 是它在 $\mathcal O=\mathcal O_{S,s}$ 上定义的\Term{group-scheme}{群概形}，则\Term{group-scheme}{群概形} $A^\natural_{\mathcal O}\otimes_{\mathcal O}\widehat{\mathcal O}$ 连同其\Term{extension}{扩张}结构，\Term{canonical}{典范}地\Term{isomorphism}{同构}于前面引入的\Term{raynaud-extension}{Raynaud 扩张} $A^\natural_{\widehat{\mathcal O}}$。
\end{SGAresult}
\begin{proof}
由 7.2.1，后一事实立即从第一个断言得出，因为根据 (13.3.6)，$A^\natural_{\mathcal O}\otimes_{\mathcal O}\widehat{\mathcal O}$ 沿\Term{special-fiber}{特殊纤维}的\Term{formal-completion}{形式完备化}为 $\widehat A_{\mathcal O}$。为了证明 $A^\natural$ 在 $s$ 附近的\Term{algebra}{代数}性，仿照 7.2.1，立即归结为证明 $A^\natural$ 的\Term{abelian-quotient}{阿贝尔商} $B$ 的\Term{algebra}{代数}性；而这将由下面的\Term{lemma}{引理}得出。
\end{proof}

\begin{SGAresult}{\Term{lemma}{引理}}{13.5.2}
设 $B$ 是\Term{analytic-space}{解析空间} $S$ 上的一个\Term{abeloid-analytic-group}{类阿贝尔解析群}，$s$ 为 $S$ 的一点，且 $B$ 沿\Term{fiber}{纤维} $B_s$ 的\Term{formal-completion}{形式完备化} $\widehat B$ \Term{algebraizable}{可代数化}。这意味着 $B_n=B\times_S S_n$ \Term{algebraizable}{可代数化}，且将 $\widehat B$ 看成 $\widehat{\mathcal O}_{S,s}$ 上的\Term{formal-scheme}{形式概形}时，它\Term{algebraizable}{可代数化}，即来自 $\widehat{\mathcal O}$ 上的一个\Term{abelian-scheme}{阿贝尔概形} $B_{\widehat{\mathcal O}}$。再假设 $B_{\widehat{\mathcal O}}$ 在 $\widehat{\mathcal O}$ 上\Term{projective}{射影}。则 $B$ 在 $s$ 附近相对于 $S$ \Term{projective}{射影}。
\end{SGAresult}
\begin{proof}
假设意味着存在\Term{formal-analytic-space}{形式解析空间} $\widehat B$ 的一个\Term{polarization}{极化}，即一列可以“\Term{gluing}{粘合}”的 $B_n$ 的\Term{polarization}{极化} $p_n$。结论意味着在 $s$ 的一个\Term{neighborhood}{邻域}上存在 $B$ 的一个\Term{polarization}{极化} $p$。设 $P$ 为 $B$ 在 $S$ 上的\Term{relative-polarization}{相对极化}构成的 $S$ 上\Term{analytic-space}{解析空间}，
\SourcePage{511}
它在 $S$ 上\Term{analytic-space}{解析空间} $S'$ 的\Term{category}{范畴}上\Term{represent}{表示}\Term{functor}{函子}
\[
 S'\longmapsto\{B_{S'}\text{ 在 }S'\text{ 上的极化}\}.
\]
这个\Term{functor}{函子}是 $S$ 的 Néron--Severi \Term{functor}{函子} $\operatorname{NS}_{B/S}$ 的一个\Term{open-subfunctor}{开子函子}。最简单的标准论证——这里不再展开——表明，$\operatorname{NS}_{B/S}$，因而 $P$，可由 $S$ 上的一个\Term{net}{净}（即\Term{unramified}{非分歧}）\Term{analytic-space}{解析空间}\Term{represent}{表示}。因此，假设就是说 $P$ 在 $s$ 处有一个 $S$ 上的\Term{formal-section}{形式截面} $p$。因为 $P$ 在 $S$ 上\Term{net}{净}，因而\Term{locally-quasi-finite}{局部拟有限}，由 [2, 第 18 讲，第 1 节]，$P$ 在 $s$ 处的任一 $S$ 上\Term{formal-section}{形式截面}都来自 $s$ 附近的一个\Term{analytic-section}{解析截面}，即 $B$ 在 $s$ 的一个\Term{neighborhood}{邻域}上的\Term{polarization}{极化}；这就得到了所述结论。
\end{proof}

\textbf{注记 13.5.3。}
\textbf{a)} 若在 13.5.2 中去掉 $B_{\widehat{\mathcal O}}$ 在 $S$ 上\Term{projective}{射影}的假设，仍可断言 $B$ 在 $s$ 附近\Term{algebraizable}{可代数化}。为此，使用一个未发表的结果（由 M.~Artin 和 P.~Deligne 各自独立地利用第 VI 讲第 2 节的 $\underline T^i$ \Term{sheaf}{层}理论得到）：若 $X$ 在 $S$ 上\Term{proper}{固有}，且 $X_{\mathrm{red}}$ \Term{algebraizable}{可代数化}，则 $X$ 也\Term{algebraizable}{可代数化}。于是可以归结到 $S$ 为\Term{reduced}{约化}的情形。另一方面，当 $S$ \Term{normal}{正规}，因而 $\widehat{\mathcal O}$ \Term{normal}{正规}时，已知 $\widehat{\mathcal O}$ 上的任一\Term{abelian-scheme}{阿贝尔概形}在 $\widehat{\mathcal O}$ 上\Term{projective}{射影}（例如参见 [23]），故可应用 13.5.2。为了得到任意\Term{reduced}{约化} $S$ 的情形，从 $S$ 的\Term{normalization}{正规化} $S'$ 出发作非\Term{flat}{平坦}的\Term{descent}{下降}，所依据的是下面这个较为微妙的结果，我们在此不予证明：\Term{noetherian-scheme}{诺特概形}之间的一个\Term{proper-morphism}{固有态射} $f:S'\to S$，若为\Term{epimorphism}{满态射}（即 $\mathcal O_S\to f_*(\mathcal O_{S'})$ 为\Term{injective}{单射}），
\SourcePage{512}
则它是\Term{relative-abelian-scheme}{相对阿贝尔概形}的\Term{fibered-category}{纤维范畴}的\Term{effective-descent-morphism}{有效下降态射}。

\textbf{b)} 看来，13.5.2 的证明方法应当也能对 $S$ 上任一\Term{proper}{固有}\Term{analytic-space}{解析空间} $B$ 证明类似结果。无疑更微妙的问题是：若 $X$ 是 $S$ 上的\Term{proper}{固有}\Term{analytic-space}{解析空间}，且 $X$ 沿 $X_s$ 的\Term{formal-completion}{形式完备化} $\widehat X$ \Term{algebraizable}{可代数化}，那么 $X$ 是否在 $s$ 附近相对于 $S$ \Term{algebraizable}{可代数化}？这里只指出，$X$ 在 $S$ 上的\Term{relative-algebraizability}{相对可代数化}概念，应取比 [10] 中更广的意义：即在 $s$ 附近，$X$ 来自 M.~Artin [1] 意义下的一个 $\mathcal O_{S,s}$ 上的\Term{algebraic-space}{代数空间}。

\textbf{13.5.4。}置身于 13.5.1 的条件之下。另对每个整数 $n>0$，选取 $\mathbb C$ 上群的标准\Term{isomorphism}{同构}
\[
 \mu_{n\mathbb C}\xleftarrow{\sim}(\mathbb Z/n\mathbb Z)_{\mathbb C},
 \qquad m\longmapsto\exp(2\mathrm i\pi m/n).
\]
对固定数 $\ell$ 的各个幂取\Term{limit}{极限}，便得到 $\mathbb C$ 上 $\ell$ 进\Term{sheaf}{层}的\Term{canonical}{典范}\Term{isomorphism}{同构}
\[
 \mathbb Z_\ell(1)_{\mathbb C}\xrightarrow{\sim}\mathbb Z_{\ell\mathbb C},
\]
以下用它来认同两端。

对\Term{exact-sequence}{正合列} (13.1.3) 应用 $\underline{\operatorname{Ext}}^i_{\mathbb Z}(\mathbb Z/n\mathbb Z,-)$ 的\Term{exact-sequence}{正合列}，得到一个\Term{canonical}{典范}\Term{isomorphism}{同构}
\SourcePage{513}
\begin{equation}
 {}_nA^0\xrightarrow{\sim}\underline R\otimes_{\mathbb Z}\mathbb Z/n\mathbb Z.
 \tag{13.5.5}
\end{equation}
当 $n$ 遍历给定\Term{prime-number}{素数} $\ell$ 的各个幂时，便得到 $S$ 上群的\Term{projective-system}{射影系统}的\Term{isomorphism}{同构}
\begin{equation}
 T_\ell(A^0)\xrightarrow{\sim}\underline R\otimes_{\mathbb Z}\mathbb Z_\ell,
 \tag{13.5.6}
\end{equation}
其中将 $\mathbb Z_\ell$ 看成 $\mathbb Z/\ell^{n+1}\mathbb Z$ 的\Term{projective-system}{射影系统}。在这里所考察的情形，$A$ 来自 $S$ 上的一个\Term{separated}{分离}\Term{group-scheme}{群概形}，于是立即可见，各 $\underline R\otimes_{\mathbb Z}\mathbb Z/n\mathbb Z$ 可看成\Term{relative-scheme}{相对概形}，从而\Term{signature}{符号} $(\underline R\otimes_{\mathbb Z}\mathbb Z_\ell)_{\mathcal O}$ 有了意义（而\Term{generic}{一般}来说 $R_{\mathcal O}$ 并没有意义！）。由此得到一个\Term{canonical}{典范}\Term{isomorphism}{同构}
\begin{equation}
 T_\ell(A_{\mathcal O})\simeq
 (\underline R\otimes_{\mathbb Z}\mathbb Z_\ell)_{\mathcal O}.
 \tag{13.5.7}
\end{equation}
由第 2 节及 13.3 的构造立即得出，上述\Term{isomorphism}{同构}将\Term{fixed-part}{固定部分}映到\Term{fixed-part}{固定部分}，将\Term{toric-part}{环面部分}映到\Term{toric-part}{环面部分}：
\begin{equation}
 \begin{aligned}
 T_\ell(A_{\mathcal O})^f&\simeq
 (\underline R^f\otimes_{\mathbb Z}\mathbb Z_\ell)_{\mathcal O},\\
 T_\ell(A_{\mathcal O})^t&\simeq
 (\underline R^t\otimes_{\mathbb Z}\mathbb Z_\ell)_{\mathcal O}.
 \end{aligned}
 \tag{13.5.8}
\end{equation}
若 $A'$ 是满足与 $A$ 相同假设的 $S$ 上\Term{relative-analytic-group}{相对解析群}，$W_{\mathcal O}$ 是 $(A_{\mathcal O},A'_{\mathcal O})$ 被 $\mathbb G_{m\mathcal O}$ 的一个\Term{biextension}{双扩张}，因而在 $s$ 的某个\Term{neighborhood}{邻域}上定义 $(A,A')$ 被 $\mathbb G_{mS}$ 的一个\Term{biextension}{双扩张}（此时使用复解析意义下的概念，即在 $S$ 上\Term{analytic-space}{解析空间}的\Term{site}{景}中，对某种适当的\Term{topology}{拓扑}进行讨论……），则根据第 VIII 讲 2.2，$W$ 定义\Term{pairing}{配对}
\SourcePage{514}
\begin{equation}
 (\underline R\otimes\mathbb Z_\ell)\times
 (\underline R'\otimes\mathbb Z_\ell)
 \simeq T_\ell(A)\times T_\ell(A')
 \longrightarrow\mathbb Z_\ell(1)_S\simeq\mathbb Z_{\ell S}.
 \tag{13.5.9}
\end{equation}
同样显然，相应的\Term{pairing}{配对}
\begin{equation}
 T_\ell(A_{\mathcal O})\times T_\ell(A'_{\mathcal O})
 \longrightarrow\mathbb Z_\ell(1)_S
 \tag{13.5.10}
\end{equation}
正是上述文献中与 $(A_{\mathcal O},A'_{\mathcal O})$ 的\Term{biextension}{双扩张} $W_{\mathcal O}$ 关联的\Term{pairing}{配对}。此外，不难看出，只须\Term{restriction}{限制}在 $s$ 的某个\Term{neighborhood}{邻域}内，各\Term{pairing}{配对} (13.5.9) 就都来自一个唯一的\Term{pairing}{配对}（它不再依赖 $\ell$）
\begin{equation}
 \underline R\otimes\underline R'\longrightarrow\mathbb Z.
 \tag{13.5.11}
\end{equation}

具备这些准备后，便可利用关于 $T_\ell(A_{\mathcal O})$ 的\Term{algebra}{代数}结果，得到关于\Term{sheaf}{层} $\underline R$ 的结论。由此可得，\Term{pairing}{配对} (13.5.11) 在 $\underline R^t\times\underline R'^f$ 和 $\underline R^f\times\underline R'^t$ 上为零（7.4.8）；当 $S$ \Term{nonsingular}{非奇异}且\Term{dimension}{维数}为 $1$，并且在 $U=S-\{s\}$ 上，$A,A'$ 是由 $W$ 使之相互\Term{dual}{对偶}的两个\Term{relative-abelian-scheme}{相对阿贝尔概形}时，则在 $s$ 的足够小\Term{neighborhood}{邻域}上，$\underline R^t_U$ 与 $\underline R'^f_U$，以及 $\underline R^f_U$ 与 $\underline R'^t_U$，恰好互为\Term{annihilator}{零化子}（7.4.8）。因为 (13.5.11) 在 $U$ 上为\Term{perfect-duality}{完美对偶}，于是得到 $U$ 上 $\underline R/\underline R^t$ 与 $\underline R'^f$ 之间的\Term{duality}{对偶}。特别地，这蕴涵 $\underline R/\underline R^t$ 为\Term{constant}{常值}的，更进而蕴涵 $\pi_1(U,t)$ 在 $\underline R_t$ 上的\Term{action}{作用}为二阶\Term{unipotent}{幺幂}。同样可见，在 $U$ 上，$\underline R/\underline R^f$ 与 $\underline R'^t$ 之间有\Term{perfect-duality}{完美对偶}；后者又\Term{isomorphism}{同构}于 $S$ 上的\Term{lattice}{格} $M_S$ 的\Term{dual-torus}{对偶环面}，其中 $M$ 是 (5.5.4) 所定义的群，因此得到\Term{canonical}{典范}\Term{isomorphism}{同构}
\SourcePage{515}
\begin{equation}
 \underline R_U/\underline R^f_U\simeq M_U.
 \tag{13.5.12}
\end{equation}
而 (13.3.4) 的\Term{homomorphism}{同态} $i$ 可以写成\Term{homomorphism}{同态}
\begin{equation}
 i:M_U\longrightarrow A^\natural;
 \tag{13.5.13}
\end{equation}
下一节将在一个稍有不同的背景中再次遇到它。
